\Needspace{10\baselineskip}
\section{完成情况与验收边界}
\label{sec:status}
本章按三态口径盘点完成度：已实现并验证、已实现未验证、未实现。
口径以主计划（2026-10-28）与三态盘点（2026-10-26）为准。

\subsection{完成度矩阵}
完成度矩阵见表~\ref{tab:status}。
\begin{table}[H]\centering\small
\caption{完成度矩阵（三态口径，含口径日期）}\label{tab:status}
\begin{tabularx}{\linewidth}{@{}P{14mm}Y@{}}\toprule
状态 & 事项\\\midrule
已实现并验证 & 核与最小 SoC 收口（2026-10-07）；视觉 RTL 与真实视频链路；视觉 RTL tb、视觉 Python 与 INT8 参考测试全量回归；FILE/LIVE 原型与前端展示离板验证；网页端真机联调（2026-10-30 联调、2026-11-03 工单下发与 CSV 导出复测）；板上硬件 CNN 四例黄金与 10{,}000 次寄存器读取（2026-10-26/27）；多目标 12 轮 85 结果与整数参考一致（2026-10-28）；CNN r6 训练与校准 33/33（2026-10-28）\\
已实现未验证 & 含 DMA 边沿修复的早期位流与 trained 包（构建完成，当时未上板；多目标修复已随 2026-10-28 共存位流上板）；正式工位构图与拍摄条件；网页端连续多轮稳定性（2026-10-30 同会话第 5 轮出现任务错误、随后 DMA 超时，根因未定位）\\
未实现 & 全链整机验收；正式 INT8 模型与生产黄金；换实物泛化验收；工位实测矩阵（N10--N12）；连续运动的触发、关联与去重\\\bottomrule
\end{tabularx}\end{table}

\subsection{结果讨论}
核与协处理器的正确性证据链完整：架构测试签名、CoreMark 交叉核对、量化逐位黄金、板上黄金对照逐级覆盖。
推理正确性与识别精度是两件事：黄金对照证明整数链与硬件输出一致，不证明识别精度；留出图对照给出诚实的精度快照（8/9），正式精度报告待真实数据与正式工位。
工单判定与记录链路经离板与软件联调验证；网页端已接入真实视频流、板端服务与记录导出，并完成页面工单下发；
连续多轮稳定性未验收，故整机验收不能宣称完成。

\subsection{验收边界声明}
本报告区分四级证据：离板仿真、软件联调、单目标板上推理、整机验收。
PROVISIONAL 黄金对拍不等同于识别精度，也不等同于整机验收。
性能指标均为实测留档口径，外推值单独标注。
\FloatBarrier
